\scanpage{89}
\mspage{12}
$\widehat\Pi_1(C,I')$ est déterminée \inserted{via (39) ou (40)}, car \unsure{l' pour|si pour} $x,y\in C$, l'application
\[
\operatorname{Hom}_{\widehat\Pi_1(C)}(x,y)\longrightarrow\operatorname{Hom}_{\widehat\Pi_1(U_{0,3})}(x,y)
\]
est injective!

\unsure{J'avais|J'ai} \unsure{pas} observé, \unsure{car|que} le saturé $I'$ de $I$ \inserted{dans $C$} n'est pas tjrs \deleted{fini} contenu dans \deleted{$\unsure{X}$} $C_K$ --- mais à vrai dire, rien ne m'obligeait à remplacer $I$ par $I'$ pour faire mon raisonnement: il reste vrai que la connaissance de l'action de $\Gamma$ sur sur $\widehat\Pi_1(U_{0,3},I_0)$ implique la connaissance de son action sur $\widehat\Pi_1(X,I)$. Mais il est vrai sûrement que l'étude des relations entre $\widehat\Pi_1(C;\deleted{\unsure{x}}I)$ et $\widehat\Pi_1(U_{0,3},I_0)$ en tant que groupoïdes à groupe d'opérations $\Gamma$, est facilité grandement en remplaçant $I$ par son saturé $I'$, $=$ si celui-ci n'est pas $\subset\deleted{X}C_K$. (Au besoin, si on est ennuyé par des sommets qui se balladent sous l'action de $\Gamma$, il n'y a
\margin{88}
